\documentclass[11pt]{article}
\usepackage[margin=1in]{geometry}
\usepackage{amsmath,amssymb}
\usepackage{booktabs}
\usepackage{hyperref}
\usepackage{xcolor}
\usepackage{enumitem}
\usepackage{microtype}
\usepackage{titlesec}

\titleformat{\section}{\normalfont\Large\bfseries}{\thesection}{1em}{}
\titleformat{\subsection}{\normalfont\large\bfseries}{\thesubsection}{1em}{}

\title{\textbf{Quantum-ICL: Simulator-Verified In-Context Learning and Lightweight Fine-Tuning for Quantum Circuit Synthesis}}
\author{Ivan Ge \\ \small Stanford University \\ \small \texttt{ivange@stanford.edu}}
\date{2026}

\hypersetup{colorlinks=true, urlcolor=blue, linkcolor=black, citecolor=black}

\begin{document}
\maketitle

\begin{abstract}
\noindent
We ask whether a \emph{frozen} large language model can learn to synthesize quantum circuits through a closed feedback loop with an exact quantum simulator -- no policy gradients, no preference data, just verified examples reused as in-context exemplars. We build a controlled benchmark that varies model capability (gpt-4o-mini $\to$ Gemini-3-Flash $\to$ Qwen2.5-7B) and task difficulty (stabilizer-state preparation, 1-qubit Clifford, 1-qubit Clifford+T) under a clean $2\times2$ ablation of two ICL mechanisms -- \textbf{verifier-feedback self-refinement} and \textbf{structural retrieval} from a growing verified library. We then compare ICL to lightweight LoRA supervised fine-tuning on simulator-verified synthetic data, scaling the training set from 0 to 3600 examples. Across 100 tasks per cell $\times$ 12 (condition, tier) cells $\times$ 5 backbones, three findings emerge: (i) verifier-feedback self-refinement is the most consistent lever; (ii) structural retrieval is \emph{useless or harmful in isolation} but stacks strongly with feedback on hard tiers; (iii) SFT on verified synthetic data is \textbf{non-monotone} -- 300 examples is catastrophic, $\geq900$ stabilizes, and 1800--3600 yields the strongest combined ICL+SFT performance. An \textsc{OracleRetrieval} probe shows our target-computable features are within $\pm1$ task of perfect oracle features: the retrieval bottleneck is not feature engineering but library size and exemplar transferability. Total compute: $\approx\$3$ of API spend on the main sweeps and $\approx25$ GPU-hours on a single A100 80\,GB.
\end{abstract}

\section{Introduction}

Quantum circuit synthesis -- given a target state or unitary, produce a circuit over an allowed gate set that realizes it -- is a foundational problem normally attacked with hand-crafted compilation (Gottesman--Knill tableau reduction for Cliffords, Solovay--Kitaev decomposition for Clifford+T, modern superoptimizers like BQSKit). It has three properties that make it an unusually good testbed for the broader question of \emph{whether frozen LLMs can self-improve through tool-mediated verification}:
\begin{enumerate}[leftmargin=*,topsep=2pt]
    \item \textbf{Cheap, perfect verification.} Exact simulation ($\leq 4$ qubits) gives binary success and a graded fidelity signal at sub-second cost.
    \item \textbf{Compact, well-typed solutions.} Circuits are short structured JSON, friendly to LLM emission and easy to parse strictly.
    \item \textbf{A natural difficulty knob.} Gate set, qubit count, and circuit depth interpolate from trivial (graph states) to intractable in a tunable way.
\end{enumerate}

We use this testbed to isolate four mechanisms that have been proposed for closed-loop frozen-LLM improvement and to ask, for each, \emph{under what regime does it actually help?}

\paragraph{Contributions.}
\begin{description}[leftmargin=*,topsep=2pt,itemsep=0pt]
    \item[\textbf{C1.}] A modular, fully reproducible benchmark (12-module Python package, 28 unit tests) implementing four difficulty tiers, four orthogonal experimental conditions (a $2\times2$ ablation of feedback $\times$ retrieval), three retrieval strategies, four LLM backends, and incremental persistence so multi-hour runs survive interruption.
    \item[\textbf{C2.}] A 100-task confirmatory sweep across three commercial/open models, demonstrating that \textbf{verifier feedback is the dominant ICL lever} while structural retrieval is regime-dependent.
    \item[\textbf{C3.}] A diagnostic showing that target-only structural features $\approx$ hidden-generator oracle features for retrieval -- the bottleneck is exemplar transfer, not feature engineering.
    \item[\textbf{C4.}] A 5-point SFT-data-size sweep ($0/300/900/1800/3600$ verified synthetic examples) on Qwen2.5-7B + LoRA, revealing a non-monotone interaction between SFT and ICL.
\end{description}

\section{Method}

\subsection{Task tiers}

We define four task tiers. Each tier specifies (i) a gate set, (ii) a hidden generator distribution, (iii) a target representation shown to the LLM, and (iv) a verification predicate.

\begin{center}
\small
\begin{tabular}{lllll}
\toprule
\textbf{Tier} & \textbf{Gate set} & \textbf{Generator} & \textbf{Target} & \textbf{Verify} \\
\midrule
A & H, CZ & Graph + canonical prep & Graph adjacency & State fidelity \\
\textbf{B} & H, S, CX, CZ & Random circuit, len 2--6 & State-vector amplitudes & State fidelity \\
\textbf{C-lite} & H, S, X, Y, Z & 1-qubit, len 1--6 & $2\times2$ unitary & Process fidelity \\
\textbf{D-lite} & H, S, T, X, Y, Z & 1-qubit, len 1--6 & $2\times2$ unitary & Process fidelity \\
\textbf{D-mid} & H, S, T, CX, CZ & 2-qubit, depth 6--10 & $4\times4$ unitary & Process fidelity \\
\bottomrule
\end{tabular}
\end{center}

Tier A is for pipeline sanity (already saturated for every commercial model) and is excluded from the main study. D-mid is implemented for the calibration appendix.

Verification uses state fidelity $|\langle\psi_{\text{target}}|\psi_{\text{out}}\rangle|$ for state tiers and global-phase-invariant Hilbert--Schmidt process fidelity $|\mathrm{Tr}(U_{\text{target}}^\dagger U_{\text{out}})|/2^n$ for unitary tiers. The success threshold is $F > 0.999$.

\subsection{The $2\times2$ ICL ablation}

We factorially cross two binary axes -- \textbf{feedback on/off} and \textbf{structural retrieval on/off}:
\begin{enumerate}[leftmargin=*,topsep=2pt,itemsep=0pt]
    \item \texttt{zero\_shot} -- empty prompt, no feedback.
    \item \texttt{feedback\_only} -- failed attempts feed back the prior circuit, its measured fidelity, and (for state tiers) the simulated state vector the proposed circuit actually produced.
    \item \texttt{structural\_retrieval\_only} -- top-$k$ verified examples retrieved from a same-tier library that grows online.
    \item \texttt{structural\_retrieval\_plus\_feedback} -- both mechanisms.
\end{enumerate}
Each task is granted up to 3 attempts. Same-seed task pools across all conditions and models give paired comparisons. Temperature is fixed at 0; max output tokens 1024.

\subsection{Structural retrieval (target-only features)}

A core methodological choice is that retrieval features must be \emph{computable from the target alone}:
\begin{itemize}[leftmargin=*,topsep=2pt,itemsep=0pt]
    \item \textbf{State tiers:} number of qubits, state sparsity, nonzero count, peak amplitude, Shannon entropy.
    \item \textbf{Unitary tiers:} number of qubits, unitary sparsity, diagonal mass $\sum_i |U_{ii}|^2/d$, Frobenius distance to identity, eigenvalue-phase vector.
\end{itemize}
Hidden-generator statistics (T-count, depth, gate count) are kept on the task object as \emph{oracle features} but are not used by the default retriever; they power only an \textsc{OracleRetrieval} diagnostic.

\subsection{Verifier-feedback self-refinement}

When an attempt fails verification, the next prompt is augmented with:
\begin{quote}\small\ttfamily
--- Your previous attempt was INCORRECT --- \\
Circuit you proposed: \{prev\_circuit\_json\} \\
It ran but matched the target with fidelity only $F$ (need $>$ 0.999). \\
Your circuit produced this state instead: [a0+b0j, a1+b1j, ...] \\
Compare it entry-by-entry with the target state above and change the gates...
\end{quote}
For unitary tiers the produced-state line is omitted (process fidelity only).

\subsection{LoRA supervised fine-tuning on verified synthetic data}

We generate $N \in \{0, 300, 900, 1800, 3600\}$ training examples (split evenly across B / C-lite / D-lite), where each example is a chat-format triple $(\textsc{system}, \textsc{user\_target}, \textsc{assistant\_generator\_circuit})$. The assistant turn is the hidden generator circuit -- verified-correct by construction.

LoRA adapters target attention $+$ MLP modules of Qwen2.5-7B-Instruct (\texttt{q, k, v, o, gate, up, down} proj). Rank $r=32$, $\alpha=64$, dropout 0.05, learning rate $10^{-4}$, cosine schedule with 5\% warmup, gradient checkpointing, bf16, 3 epochs (2 for $N=3600$). Effective batch size 16 (per-device 2 $\times$ grad-accum 8).

\paragraph{Critical design choice: assistant-only loss.} Our first attempt used standard \texttt{SFTTrainer} loss over the full token sequence. Training loss converged to 0.046 but produced mode collapse -- the model emitted a single fixed circuit \texttt{Y,Z,S,Y,Z} for every target. The boilerplate ($\sim$90\% of tokens) dominated the loss while the actual gate sequence was ignored. Switching to TRL's \texttt{assistant\_only\_loss=True} (masks system+user tokens) raised loss to 0.139 -- a desirable increase -- and the model began producing tier-appropriate per-target circuits.

\subsection{Backbones and reproducibility}

We use \textbf{gpt-4o-mini} (small commercial baseline), \textbf{Gemini-3-Flash} (\texttt{google/gemini-3-flash-preview}), \textbf{Qwen2.5-7B-Instruct} (base), and \textbf{Qwen2.5-7B + LoRA SFT} at five data sizes. DeepSeek-R1 and explicit reasoning models are deliberately excluded from main sweeps: pilot runs showed r1 saturating easy tiers but costing $\sim\$0.28$ per 20-task block at $\sim$90\,s/call, making large-scale ablations infeasible.

\textbf{Reproducibility infrastructure:} per-(condition, tier, model) seeds via SHA-256 hashing -- reproducible across Python processes regardless of \texttt{PYTHONHASHSEED}. \textbf{Incremental per-block persistence:} each block dumps records to \texttt{<run\_dir>/blocks/} when it completes; \texttt{--rebuild} reconstructs summaries after timeouts. 28 unit tests cover schema validation, verification correctness across all tiers, perturbation failures, global-phase invariance, gate-set enforcement, retrieval ranking, SFT, and OracleRetrieval.

\section{Experiments and Results}

All experiments use 100 tasks per cell (except where noted), seed 42, up to 3 attempts, temperature 0.

\subsection{Experiment 1: $2\times2$ ICL ablation $\times$ three model backbones}

Twelve cells per model; solve counts out of 100:

\begin{center}\small
\begin{tabular}{llrrr}
\toprule
\textbf{Tier} & \textbf{Condition} & \textbf{gpt-4o-mini} & \textbf{Gemini-3-Flash} & \textbf{Qwen-7B (base)} \\
\midrule
B & zero\_shot & 9 & \textbf{79} & 8 \\
  & feedback\_only & 11 & 75 & 14 \\
  & structural\_only & 6 & 79 & 11 \\
  & structural+feedback & \textbf{12} & \textbf{86} & \textbf{22} \\
\midrule
C-lite & zero\_shot & 11 & 57 & 23 \\
  & feedback\_only & \textbf{15} & 74 & \textbf{30} \\
  & structural\_only & 10 & 73 & 21 \\
  & structural+feedback & \textbf{15} & \textbf{100}$^{\dagger}$ & 27 \\
\midrule
D-lite & zero\_shot & 7 & 28 & 14 \\
  & feedback\_only & \textbf{8} & 53 & 20 \\
  & structural\_only & 3 & 54 & 15 \\
  & structural+feedback & 4 & \textbf{65} & \textbf{20} \\
\bottomrule
\end{tabular}
\end{center}
{\small $^\dagger$ early-stopped at 10/10 (block saturated).}

Cost: gpt-4o-mini \$0.22, Gemini-3-Flash \$1.33, Qwen-base 0 (local).

\subsection{Experiment 2: OracleRetrieval diagnostic}

A 30-task probe on Gemini-3-Flash compares default structural retrieval (target-only features) to \textsc{OracleRetrieval} that cheats and uses hidden generator features:

\begin{center}\small
\begin{tabular}{lrrr}
\toprule
\textbf{Tier} & \textbf{Structural+fb (target)} & \textbf{OracleRetrieval (hidden)} & $\Delta$ \\
\midrule
B & 10/10$^\dagger$ & 29/30 & $\approx 0$ \\
C-lite & 24/30 & 24/30 & 0 \\
D-lite & 18/30 & 19/30 & $+1$ \\
\bottomrule
\end{tabular}
\end{center}

The structural retriever using only target-derivable features is within $\pm1$ task of the cheat. \textbf{The retrieval bottleneck is not feature quality} -- it is library size and exemplar transferability.

\subsection{Experiment 3: SFT data-size scan on Qwen2.5-7B}

Solve counts at five SFT corpus sizes (0 = base, no fine-tune); same seeds as Experiment 1:

\begin{center}\small
\begin{tabular}{llrrrrr}
\toprule
\textbf{Tier} & \textbf{Condition} & \textbf{0} & \textbf{300} & \textbf{900} & \textbf{1800} & \textbf{3600} \\
\midrule
B & zero\_shot & 8 & 7 & 17 & 20 & 16 \\
  & feedback\_only & 14 & 11 & 19 & 22 & \textbf{27} \\
  & struct\_only & 11 & 9 & 12 & 11 & \textbf{16} \\
  & struct+fb & 22 & --- & 12 & 15 & \textbf{26} \\
\midrule
C-lite & zero\_shot & 23 & 3 & 17 & 18 & 16 \\
  & feedback\_only & 30 & 9 & 34 & \textbf{35} & 34 \\
  & struct\_only & 21 & --- & 10 & 18 & 11 \\
  & struct+fb & 27 & --- & 27 & 32 & \textbf{35} \\
\midrule
D-lite & zero\_shot & 14 & 0/10$^{-}$ & 3 & 10 & \textbf{16} \\
  & feedback\_only & 20 & 0/10$^{-}$ & 10 & 18 & \textbf{22} \\
  & struct\_only & 15 & --- & 7 & 13 & 10 \\
  & struct+fb & 20 & --- & 21 & \textbf{28} & 24 \\
\bottomrule
\end{tabular}
\end{center}
{\small $^{-}$ early-stopped at 0/10; --- = block did not complete (SFT-300 hit 6\,h walltime; 7 of 12 blocks recovered via per-block persistence).}

Compute: 5 SFT trainings totalled 1\,h 24\,min A100 time. Final training loss in $[0.10, 0.15]$ for all sizes.

\section{Findings}

\paragraph{F1: Verifier feedback is the dominant ICL lever.} Across all 36 (model, tier) pairs at 100 tasks, adding feedback to zero-shot never decreases solve rate by more than 4 points and increases it by up to 32 points (Gemini D-lite: $28\to53$). It is the only mechanism monotone in this sense.

\paragraph{F2: Structural retrieval is regime-dependent and negative in isolation.} For all three frozen models, \texttt{structural\_only} matches or under-performs \texttt{zero\_shot} on Tier B and on D-lite. The lift comes only when combined with feedback: \texttt{structural+feedback} is the maximum solve count for \textbf{9 of 9 cells} for Gemini-3-Flash and \textbf{2 of 3 tiers} for Qwen-base.

\paragraph{F3: The retrieval bottleneck is exemplar transfer, not feature engineering.} \textsc{OracleRetrieval} $\approx$ structural retrieval on Gemini-3-Flash (within $\pm1/30$). Pushing the retrieval signal further requires more verified examples, more transferable exemplars, or a different retrieval objective -- not better features.

\paragraph{F4: SFT on verified synthetic data is non-monotone and can interfere with ICL.}
\begin{itemize}[leftmargin=*,topsep=2pt,itemsep=0pt]
    \item \textbf{300 examples is catastrophic.} D-lite zero-shot $14\to0$ (every one of the first 10 fails); C-lite zero-shot $23\to3$.
    \item \textbf{900 examples is the floor for stable training} -- performance recovers to $\approx$ base.
    \item \textbf{1800--3600 yields the strongest gains} for feedback-containing conditions (B feedback: $14\to27$).
    \item \textbf{Small/medium SFT interferes with retrieval.} B struct+fb: $22$ (base) $\to 15$ (1800) $\to 26$ (3600) -- a U-shape: light SFT teaches a circuit-emission style the model defaults to, partially ignoring retrieved exemplars; large SFT recovers context-attention.
\end{itemize}

\paragraph{F5: Tier difficulty calibration is critical for revealing ICL effects.}
Tier A saturates for every commercial model; Tier B saturates for Gemini-3-Flash (79--86\% zero-shot) but is in the sweet spot for the weaker models (8--22\%). D-lite is the \emph{only} tier where structural+feedback shows a clean monotone gain on a strong frontier model (Gemini: $28\to65$, $+37$). The methodological implication: \emph{to demonstrate verified-ICL effects empirically, calibrate the task so baseline zero-shot success lies in 20--60\%.}

\section{Discussion}

\paragraph{Why does feedback dominate over retrieval?} Feedback acts on \emph{this specific} failed attempt: the model sees the exact discrepancy and can locally edit gates. Retrieval offers a related but distinct problem the model must analogize from. Local editing is easier than analogical transfer for the models tested, especially when fidelity is high but $\ne 1$ (a few-gate fix is usually enough).

\paragraph{Why does structural retrieval fail in isolation?} The retrieved circuit solves a \emph{different} target. Without feedback, the model either copies the wrong answer or finds no useful transformation. Combined with feedback, retrieval supplies \emph{form} (gate sequences that resemble correct ones) and feedback supplies \emph{correction} (this specific attempt is wrong, here's how). They are complementary.

\paragraph{Why does SFT interfere with retrieval at intermediate scale?} A LoRA adapter trained on a single response style biases the model toward emitting that style regardless of context. With few enough examples the bias overrides the in-context exemplar. With enough examples (3600 here) the model learns the \emph{task} rather than just the \emph{style}, restoring sensitivity to retrieved exemplars.

\paragraph{Capability ladder.} Across the three frozen backbones:
\begin{itemize}[leftmargin=*,topsep=2pt,itemsep=0pt]
    \item \textbf{Below the floor} (gpt-4o-mini D-lite): no mechanism helps much. The model cannot exploit any signal.
    \item \textbf{At the edge} (Qwen-7B; Gemini D-lite): every mechanism stacks. Structural+feedback is consistently best.
    \item \textbf{At saturation} (Gemini B, C-lite): all conditions reach 70--100\%. Differences are within sampling noise.
\end{itemize}
A natural conjecture: \emph{frozen-LLM improvement mechanisms are regime-dependently additive} -- sub-additive at the floor, super-additive at the edge, irrelevant at saturation.

\section{Limitations}

\begin{itemize}[leftmargin=*,topsep=2pt,itemsep=0pt]
    \item Single seed; Wilson 95\% CI roughly $\pm 9\%$ at $n=100$. Paired McNemar tests planned.
    \item Tier D / D-mid (2-qubit Clifford+T) calibration runs are queued at submission time; full results pending.
    \item Library-size confound: retrieval-condition libraries grow online; an isolated library-size sweep is in progress.
    \item Feedback ablation pending: which feedback component carries the signal (prior circuit, fidelity number, produced state vector)?
    \item Retrieval top-$k$ fixed at 3; no diversity sweep yet.
    \item SFT-300 walltime cutoff: 5 of 12 blocks missing (trend is clear).
\end{itemize}

\section{Related Work}

Self-refinement on simulator-verified targets has antecedents in code-self-debug systems (Reflexion-style verified loops). Retrieval-augmented in-context learning has a long literature (kNN-LM, in-context demonstration selection). LLM-driven quantum circuit synthesis has appeared in recent preprints; ours appears to be the first study to factorially disentangle verifier feedback $\times$ structural retrieval at this granularity, plus the SFT-data-size interaction.

\section{Conclusion}

For frozen LLMs at the edge of capability, \textbf{verifier feedback dominates and structural retrieval helps only when stacked with feedback}. For open 7B models, \textbf{LoRA SFT on simulator-verified synthetic data is a real lever} -- but its interaction with ICL is non-monotone, and the obvious ``more data is better'' intuition fails between 0 and 900 examples. The cleanest gains come from combining feedback, retrieval, and SFT on tasks calibrated to baseline solve rates of 20--60\%. The framework, datasets, and full experimental record are reproducible end-to-end on a single A100 80\,GB for under \$20 of API spend.

\paragraph{AI \& tool disclosure.} The implementation, experiment design, and writeup were co-developed with the Claude Code agent (Claude Opus 4.7), with the author making all research-direction decisions, model and tier choices, budget calls, and final analytic interpretations. The agent contributed to code authoring, batch-job orchestration, statistical aggregation, and prose drafting; all results, plots, and runs presented are from re-executable scripts in the public repository.

\paragraph{Code and data.} \url{https://github.com/ivantqge/quantum_icl}

\end{document}
